\documentclass[11pt,a4paper]{article}
\usepackage[a4paper,margin=0.5in]{geometry}
\usepackage[utf8]{inputenc}
\usepackage{mdwlist}
\usepackage[default]{lato}
\usepackage[T1]{fontenc}
\usepackage{textcomp}
\usepackage{fontawesome}
\usepackage{enumitem}
\usepackage{newunicodechar}
\newunicodechar{（}{(}
\newunicodechar{）}{)}
\DeclareFontFamily{U}{fontawesomeOne}{}
\DeclareFontShape{U}{fontawesomeOne}{m}{n}
{<-> FontAwesome--fontawesomeone}{}
\DeclareRobustCommand\FAone{\fontencoding{U}\fontfamily{fontawesomeOne}\selectfont}
\pagestyle{empty}
\setlength{\tabcolsep}{0em}

% indentsection style, used for sections that aren't already in lists
% that need indentation to the level of all text in the document
\newenvironment{indentsection}[1]%
{\begin{list}{}%
{\setlength{\leftmargin}{#1}}%
     \item[]
}
{\end{list}}

% opposite of above; bump a section back toward the left margin
\newenvironment{unindentsection}[1]%
{\begin{list}{}%
{\setlength{\leftmargin}{-0.5#1}}%
\item[]%
}
{\end{list}}

% format two pieces of text, one left aligned and one right aligned
\newcommand{\headerrow}[2]
{\begin{tabular*}{\linewidth}{l@{\extracolsep{\fill}}r}
#1 &
#2 \\
\end{tabular*}}

% make "C++" look pretty when used in text by touching up the plus signs
\newcommand{\CPP}
{C\nolinebreak[4]\hspace{-.05em}\raisebox{.22ex}{\footnotesize\bf ++}}

% and the actual content starts here
\begin{document}

\begin{center}
	{\LARGE \textbf{FENG-AN (Alan) HSIEH}}\\
	Zhubei, Hsinchu, Taiwan
	\vspace{0.05cm}
	\\
	\raisebox{-0.2\height}{\Large \faPhoneSquare} \ \ (+886) 986161907 
    \raisebox{-0.2\height}{\Large \faEnvelopeSquare} \ \ xiefengan013@gmail.com 
    \raisebox{-0.2\height}{\Large \faLinkedinSquare} \ \ linkedin.com/in/feng-an-hsieh-a0a69a15a/
\end{center}

\subsection*{\large Summary}
\hrule
\vspace{0.3em}

\hyphenpenalty=1000
\begin{itemize}[leftmargin=1em,noitemsep]
	\item[]
    ESL Modeling Engineer with 4 years of experience delivering C++/SystemC and TLM-2.0 virtual platforms for enterprise and client SSD controllers. End-to-end owner of 6 peripheral and controller IP models delivered to firmware teams and external customers. Pragmatic automation builder using Python and LLM-assisted workflows to accelerate verification, regression, and model delivery.
\end{itemize}

% ======================== Experience Section ==========================
\subsection*{\large Experience}
\hrule
\vspace{0.3em}

\renewcommand\labelitemi{}
\renewcommand\labelitemii{$\bullet$}
\begin{itemize}[leftmargin=1em]
	\parskip=0.1em
		
	\item
	      \headerrow
	      {\textbf{SiliconMotion} -- \textbf{ESL Modeling Team}}
	      {\textbf{Zhubei, Hsinchu}}
	      \headerrow
	      {\emph{Chief Engineer}}
	      {\emph{July 2026 -- Present}}
	      \begin{itemize*}
	      	\item Architect and maintain 6 C++/SystemC TLM IP models (4 enterprise, 2 client), including the shared inter-engine messaging IP and core NAND Flash Controller (arbitration, DMA, command decode).
	      	\item Profiled the messaging IP model with VDK top-down profiling and Valgrind, and replaced per-clock-cycle \texttt{wait()} polling with coarser time-based waits, reducing scheduler wake-ups and speeding up simulation.
	      	\item Deliver verified virtual platform releases to 5+ firmware engineers and external customers; mentor engineers on ESL design patterns and TLM-2.0 modeling standards.
	      	\item Built agentic developer workflows (Claude Code, Cursor) for model boilerplate generation, spec-to-model scaffolding, and automated regression, reducing new NAND model delivery from $\sim$2 weeks to 2--3 days.
	      \end{itemize*}
	      \headerrow
	      {\emph{Senior Engineer}}
	      {\emph{November 2022 -- June 2026}}
	      \begin{itemize*}
	      	\item Built full-system virtual platforms using Synopsys VDK with processor fast models and peripheral TLMs, enabling firmware teams to bring up software months ahead of silicon tape-out.
	      	\item Conducted hardware-software co-simulation with DV test suites; verified CModels against RTL specs to isolate and resolve complex boundary bugs across HW and FW.
	      	\item Automated RTL waveform parsing (FSDB/VCD) with Python and shell scripts, speeding up test creation by 3x and saving $\sim$1 week of manual effort per cycle.
	      	\item Built internal Python utilities for platform trace analysis and debugging, accelerating issue isolation for firmware teams.
	      \end{itemize*}
\end{itemize}

% ============================ Skill Section ===============================
\subsection*{\large Skills}
\hrule
\vspace{0.3em}

\hyphenpenalty=1000
\begin{itemize}[leftmargin=1em,noitemsep]
	\item \textbf{Languages \& Modeling:}
	      C/C++, SystemC, TLM-2.0, Python, Shell, Verilog, Assembly
	\item \textbf{EDA \& Virtual Platforms:}
	      Synopsys VDK, Virtual Prototyping, HW/SW Co-simulation, Git
	\item \textbf{Protocols \& Architecture:}
	      AMBA spec \& bus integration (AXI, AHB, APB, CHI), NAND Flash Controller Architecture, DMA
	\item \textbf{Engineering Tooling:}
	      Claude Code, Cursor, Automated Test Generation, Waveform Parsing (FSDB/VCD)
    \item \textbf{Languages:}
          Mandarin (Native), English (TOEIC 780)
\end{itemize}

% ======================== Education Section ==========================
\subsection*{\large Education}
\hrule
\vspace{0.3em}

\begin{itemize}[leftmargin=1em]
	\parskip=0.1em
		
	\item
	      \headerrow
	      {\textbf{National Chung Cheng University}}
	      {\textbf{Chiayi, Taiwan}}
	      \headerrow
	      {\emph{M.E. in Electrical Engineering}}
	      {\emph{August 2020 -- July 2022}}
          \headerrow
          {\emph{Thesis: An End-To-End System For Road Agent Behavior Classification (Deep Learning, Computer Vision)}}
          {\emph{}}
          
	      \headerrow
          {\textbf{National University of Kaohsiung}}
	      {\textbf{Kaohsiung, Taiwan}}
          \headerrow
	      {\emph{B.E. in Computer Science and Information Engineering}}
	      {\emph{September 2016 -- June 2020}}
\end{itemize}

\end{document}
